\begin{statementbox}
	\textbf{\textcolor{CStatement}{条件}}
	\begin{description}[style=nextline, leftmargin=2.8em, labelsep=0.5em, itemsep=0.35em]
		\item[\textbf{\textcolor{CStatement}{条件1（函数与导数形式）：}}]
			设\(f(x)=ax^3+bx^2+cx+d\;(a\neq0)\)，其导函数\(f'(x)=3ax^2+2bx+c\)，判别式\(\Delta'=4(b^2-3ac)\)。
	\end{description}
	\textbf{\textcolor{CStatement}{结论}}
	\begin{description}[style=nextline, leftmargin=2.8em, labelsep=0.5em, itemsep=0.45em]
		\item[\textbf{\textcolor{CStatement}{结论一（零点个数分类判定）：}}]
			\textbf{\textcolor{CStatement}{直观描述：}}
			判别式控制单调，极值乘积决定零点个数。 若\(\Delta'\leq0\)，则\(f(x)\)在\(\mathbb{R}\)上单调，零点个数为1。 若\(\Delta'>0\)，设\(f'(x)=0\)的两实根\(x_1<x_2\)，则零点个数由\(f(x_1)f(x_2)\)的符号决定：
			\[
				f(x_1)f(x_2)
				\begin{cases}
					<0 \rightarrow 3,\\ =0 \rightarrow 2,\\ >0 \rightarrow 1.
			\end{cases}
			\]
			其中数字表示零点个数；当乘积为零时有两个零点（包含一个重根）。
	\end{description}

\end{statementbox}
